\documentclass[12pt]{article}
\RequirePackage{style_for_analysis}

\begin{document}


\title{Notes on Multivariable Calculus}
\author{}
\date{}
\maketitle

\section{Multivariable Differential Calculus}

\subsection*{Motivation}
Let $f: \mathbb{R}^n \to \mathbb{R}^m$ be defined by $f(\mathbf{x})=(f_1(\mathbf{x}), f_2(\mathbf{x}), \dots, f_m(\mathbf{x}))$, where each component $f_i: \mathbb{R}^n \to \mathbb{R}$ is a real-valued function. 

Regardless of whether $f$ is linear or non-linear, a natural question arises: can it be locally approximated by a linear function? And if so, how accurate is this approximation?

Ultimately, given a point $\mathbf{x}_0 \in \mathbb{R}^n$, our goal is to approximate $f$ near $\mathbf{x}_0$ using its linearization:
$$f(\mathbf{x}) \approx f(\mathbf{x}_0) + Df(\mathbf{x}_0) (\mathbf{x} - \mathbf{x}_0)$$


\begin{definition}[Differentiability]
Let $E \subseteq \mathbb{R}^n$, $f: E \to \mathbb{R}^m$ be a function, and $\mathbf{x}_0$ be an interior point of $E$. Let $L: \mathbb{R}^n \to \mathbb{R}^m$ be a linear transformation. We say $f$ is differentiable at $\mathbf{x}_0$ with derivative $L$ if 
$$\lim_{\mathbf{x} \to \mathbf{x}_0} \frac{\|f(\mathbf{x}) - f(\mathbf{x}_0) - L(\mathbf{x}-\mathbf{x}_0)\|}{\|\mathbf{x}-\mathbf{x}_0\|} = 0.$$
Equivalently,
$$\lim_{\mathbf{h}\to\mathbf{0}}\frac{\| f(\mathbf{x}_0+\mathbf{h})-f(\mathbf{0})-L(\mathbf{h}) \|}{\|\mathbf{h}\|}=0.$$
\label{differentiablity}
\end{definition}

Therefore, multivariable differentiability is defined such that the relative error between the function and its linear approximation vanishes as the distance approaches zero.

\begin{theorem}
Let $E \subseteq \mathbb{R}$, $f: E \to \mathbb{R}$ be a function, $L \in \mathbb{R}$, and $x_0$ be an interior point of $E$. TFAE:
\begin{enumerate}
    \item $f$ is differentiable at $x_0$ and $f'(x_0) = L$.
    \item $\lim_{x \to x_0} \frac{|f(x) - f(x_0) - L(x - x_0)|}{|x - x_0|} = 0$.
\end{enumerate}
\end{theorem}

\begin{dbox}
\begin{pf}
For $x\neq x_0$, write
\[\frac{f(x) - f(x_0) }{x - x_0}= L + \frac{f(x) - f(x_0) - L(x - x_0)}{x - x_0}.\]
Rearrange $L$ to the left, take the absolute value and limit $x\to x_0$, and we can show the limit on the both sides imply each other.
The second expression is exactly the special case of definition \eqref{differentiablity}.
\end{pf}
\end{dbox}



\begin{theorem}[Uniqueness of derivatives]
Let $\mathbf{x}_0 \in E$ be an interior point of $E$. 
Suppose $L_1, L_2: \mathbb{R}^n \to \mathbb{R}^m$ are linear transformations such that $f$ is differentiable at $\mathbf{x}_0$ with derivative $L_1$ and also differentiable at $\mathbf{x}_0$ with derivative $L_2$. 
Then $L_1 = L_2$.
\end{theorem}


\begin{dbox}
\begin{pf}
For $\mathbf{h}\ne \mathbf{0}$,
\[\frac{\| L_1(\mathbf{h}) - L_2(\mathbf{h}) \|}{\|\mathbf{h}\|}\le
\frac{\| f(\mathbf{x}_0+\mathbf{h})-f(\mathbf{0})-L_1(\mathbf{h}) \|}{\|\mathbf{h}\|} +
\frac{\| f(\mathbf{x}_0+\mathbf{h})-f(\mathbf{0})-L_2(\mathbf{h}) \|}{\|\mathbf{h}\|}.
\]

By sandwich, $\lim_{\mathbf{h}\to\mathbf{0}}\frac{\| L_1(\mathbf{h}) - L_2(\mathbf{h}) \|}{\|\mathbf{h}\|}=0.$

Let $\mathbf{h}=t\mathbf{v}, t\ne 0$, we have
\begin{align*}
\lim_{\mathbf{\mathbf{v}}\to\mathbf{0}}\frac{\| L_1(t\mathbf{v}) - L_2(t\mathbf{v}) \|}{\| t\mathbf{v}\|} = 
\lim_{\mathbf{\mathbf{v}}\to\mathbf{0}}\frac{|t| \|L_1(\mathbf{v}) - L_2(t\mathbf{v})\|}{|t| \|\mathbf{v}\|} \quad (\text{$L_1$ and $L_2$ are linear.})
\end{align*}

Hence, for $\mathbf{v}\ne\mathbf{0}$, $L_1(\mathbf{v})=L_2(\mathbf{v})$, which is also true when $\mathbf{v}=\mathbf{0}.$
\end{pf}
\end{dbox}

We may prove by contradiction that $f$ is not differnetiable by the above theorem.

\begin{theorem}
If $f$ is differentiable at $\mathbf{x}_0$, then $f$ is continuous at $\mathbf{x}_0$.
\end{theorem}

\begin{dbox}
\begin{pf}
For $\mathbf{x}\neq \mathbf{x}_0$, write
\[\|f(\mathbf{x}) - f(\mathbf{x}_0)-L(\mathbf{x} - \mathbf{x}_0)\| = \underbrace{\frac{\|f(\mathbf{x}) - f(\mathbf{x}_0) - L(\mathbf{x} - \mathbf{x}_0)\|}{\|\mathbf{x} - \mathbf{x}_0\|}}_{\to 0 \text{ by assumption}} \underbrace{\|\mathbf{x} - \mathbf{x}_0\|}_{\to 0 \text{ as } \mathbf{x}\to\mathbf{x}_0}.\]
Next, use
\[\|f(\mathbf{x}) - f(\mathbf{x}_0)\|\le \underbrace{\|f(\mathbf{x}) - f(\mathbf{x}_0)-L(\mathbf{x} - \mathbf{x}_0)\|}_{\to 0 \text{ from the previous equality}} + \underbrace{\|L(\mathbf{x} - \mathbf{x}_0)\|}_{\to 0 \text{ as } \mathbf{x}\to\mathbf{x}_0}.\]
By Sandwich, $\lim_{\mathbf{x}\to\mathbf{x}_0}f(\mathbf{x})=f(\mathbf{x}_0).$
\end{pf}
\end{dbox}


\begin{definition}[Directional Derivative]
Let $E \subseteq \mathbb{R}^n$, $f: E \to \mathbb{R}^m$ be a function, and $\mathbf{x}_0$ be an interior point of $E$.
Fix $\mathbf{v}\in\mathbb{R}^n$, if the limit
$$D_\mathbf{v} f(\mathbf{x}_0) = \lim_{t \to 0} \frac{f(\mathbf{x}_0 + t\mathbf{v}) - f(\mathbf{x}_0)}{t}$$
exists, we call this limit the directional derivative of $f$ at $\mathbf{x}_0$ in the direction $\mathbf{v}.$
\end{definition}


\begin{definition}[Partial Derivative]
Let $E \subseteq \mathbb{R}^n$, $f: E \to \mathbb{R}^m$ be a function, and $\mathbf{x}_0$ be an interior point of $E$.
Let $1\le j\le n$, if the limit
$$\frac{\partial f}{\partial x_j}(\mathbf{x}_0) = \lim_{t \to 0} \frac{f(\mathbf{x}_0 + t\mathbf{e}_j) - f(\mathbf{x}_0)}{t}$$
exists, we call this limit the partial derivative of $f$ with respect to the variable $x_j$.
\end{definition}

So partial derivative is just a special case of directional derivative.

\begin{theorem}
If $f$ is differentiable at $\mathbf{x}_0$, then $f$ is differentiable in the direction $\mathbf{v}$ at $\mathbf{x}_0$, and:
$$D_\mathbf{v} f(\mathbf{x}_0) = f'(\mathbf{x}_0)\cdot\mathbf{v}.$$
\end{theorem}

\begin{dbox}
\begin{pf}
For $\mathbf{h}\ne 0,$ let $\mathbf{h}=t\mathbf{v},\, t\ne 0,$ write
\[
\frac{f(\mathbf{x}_0 + t\mathbf{v})-f(\mathbf{x}_0)}{t} 
= \frac{f(\mathbf{x}_0 + t\mathbf{v})-f(\mathbf{x}_0)-L(t\mathbf{v})}{t}
+ \frac{L(t\mathbf{v})}{t}.
\]
Rearrange $L(\mathbf{v})$ to the left, and take the Euclidean norm and $t\to 0$, the RHS converge to 0 by differentiablity ($\|\mathbf{v}\|\ne 0$). So we have
\[
\lim_{t\to 0}\frac{f(\mathbf{x}_0 + t\mathbf{v})-f(\mathbf{x}_0)}{t} = L(\mathbf{v}) = f^\prime(\mathbf{x}_0)\cdot\mathbf{v}.
\]
\end{pf}
\end{dbox}



We may also disprove the differentiability of $f$ provided its directional derivative exists.
The next question is: how to write the linear map $f'(\mathbf{x}_0)$ explicitly?

\begin{theorem}
Let $E \subseteq \mathbb{R}^n$, $f: E \to \mathbb{R}^m$ be a function, and $\mathbf{x}_0$ be an interior point of $E$.
If all the partial derivatives of $f$ exist on $E$ and are continuous at $\mathbf{x}_0$, then $f$ is differentiable at $\mathbf{x}_0$, and
$$f'(\mathbf{x}_0)\cdot\mathbf{v} = \sum_{j=1}^n v_j \frac{\partial f}{\partial x_j}(\mathbf{x}_0).$$
\end{theorem}


\begin{dbox}
\begin{pf}
\noindent \textbf{Step 1: Setup}\\
Let $L(\mathbf{v}) = \sum_{j=1}^n v_j \frac{\partial f}{\partial x_j}(\mathbf{x}_0)$ be our candidate for the derivative.
Since $\mathbf{x}_0$ is an interior point of $E$, there exists $r > 0$ such that the ball $B(\mathbf{x}_0, r) \subseteq E$.

For $\epsilon > 0$, since each partial derivative $\frac{\partial f_i}{\partial x_j}$ is continuous at $\mathbf{x}_0$, there exists $\delta_{i,j} > 0$ such that $\|\mathbf{x} - \mathbf{x}_0\| < \delta_{i,j}$ implies
$$\left| \frac{\partial f_i}{\partial x_j}(\mathbf{x}) - \frac{\partial f_i}{\partial x_j}(\mathbf{x}_0) \right| < \frac{\epsilon}{mn}$$
for all $1 \le i \le m$ and $1 \le j \le n$.

Let $\delta = \min\{r, \delta_{i,j} \mid 1 \le i \le m, 1\le j \le n\}$.
Now, fix $\mathbf{x} \in E$ such that $\|\mathbf{x} - \mathbf{x}_0\| < \delta$. Let $\mathbf{v} = \mathbf{x} - \mathbf{x}_0 = \sum_{j=1}^n v_j \mathbf{e}_j$. Note that $|v_j| \le \|\mathbf{x} - \mathbf{x}_0\| < \delta$.
We construct a path from $\mathbf{x}_0$ to $\mathbf{x}$ along the coordinate axes. Define the sequence of points:
\begin{align*}
\mathbf{x}^{(0)} &= \mathbf{x}_0 \\
\mathbf{x}^{(1)} &= \mathbf{x}_0 + v_1 \mathbf{e}_1 \\
\mathbf{x}^{(2)} &= \mathbf{x}_0 + v_1 \mathbf{e}_1 + v_2 \mathbf{e}_2 \\
&\vdots \\
\mathbf{x}^{(j)} &= \mathbf{x}^{(j-1)} + v_j \mathbf{e}_j \\
&\vdots \\
\mathbf{x}^{(n)} &= \mathbf{x}_0 + \sum_{j=1}^n v_j \mathbf{e}_j = \mathbf{x}
\end{align*}

We can express the difference $f(\mathbf{x}) - f(\mathbf{x}_0)$ as a telescoping sum:
$$f(\mathbf{x}) - f(\mathbf{x}_0) = \sum_{j=1}^n \left( f(\mathbf{x}^{(j)}) - f(\mathbf{x}^{(j-1)}) \right)$$

For each component $i$, consider $f_i(\mathbf{x}^{(j)}) - f_i(\mathbf{x}^{(j-1)})$. 
Let $\varphi_i^j(t) = f_i(\mathbf{x}^{(j-1)} + t v_j \mathbf{e}_j)$, then $f_i(\mathbf{x}^{(j)}) - f_i(\mathbf{x}^{(j-1)}) = \varphi^j_i(1)-\varphi^j_i(0)$.
By the Mean Value Theorem, $\exists t_j\in (0,1)$ such that $\varphi^j_i(1)-\varphi^j_i(0) = (\varphi^j_i)'(t_j)$. 

\vspace{0.8em}
\noindent \textbf{Step 2: Continuity of $f_i$ along the coordinate $\mathbf{x}_j$}\\
Applying Lemma \eqref{coordinate chain rule}: 
$$
f_i(\mathbf{x}^{(j)}) - f_i(\mathbf{x}^{(j-1)}) = (\varphi^j_i)'(t_j) = \frac{\partial f_i}{\partial x_j}(\boldsymbol{\xi}_j) v_j,
$$ 
where $\boldsymbol{\xi}_j = \mathbf{x}^{(j-1)} + t_j v_j \mathbf{e}_j$. 

Note that
\begin{align*}
\|\mathbf{x}^{(j-1)} + t_j v_j \mathbf{e}_j - \mathbf{x}_0\| &= \|v_1 \mathbf{e}_1 + v_2 \mathbf{e}_2 + \dots + v_{j-1} \mathbf{e}_{j-1} + t_j v_j \mathbf{e}_j\| \\
&\le \left\| \sum_{i=1}^n v_i \mathbf{e}_i \right\| \quad (\because |t_j| < 1) \\
&= \|\mathbf{x} - \mathbf{x}_0\| \le \delta,
\end{align*}
which implies
\begin{align*}
\left| f_i(\mathbf{x}^{(j)}) - f_i(\mathbf{x}^{(j-1)}) - \frac{\partial f_i}{\partial x_j}(\mathbf{x}_0) v_j \right| 
&= \left| \frac{\partial f_i}{\partial x_j}(\mathbf{x}^{(j-1)} + t_j v_j \mathbf{e}_j) - \frac{\partial f_i}{\partial x_j}(\mathbf{x}_0) \right| |v_j| \\
&< \frac{\epsilon}{mn} |v_j|.
\end{align*}

\vspace{0.8em}
\noindent \textbf{Step 3: Bound the error for each component $f_i$}\\
\begin{align*}
\left| f_i(\mathbf{x}) - f_i(\mathbf{x}_0) - \sum_{j=1}^n \frac{\partial f_i}{\partial x_j}(\mathbf{x}_0) v_j \right| &= \left| \sum_{j=1}^n \left( f_i(\mathbf{x}^{(j)}) - f_i(\mathbf{x}^{(j-1)}) - \frac{\partial f_i}{\partial x_j}(\mathbf{x}_0) v_j \right) \right| \\
&\le \sum_{j=1}^n \left| \frac{\partial f_i}{\partial x_j}(\boldsymbol{\xi}_j) - \frac{\partial f_i}{\partial x_j}(\mathbf{x}_0) \right| |v_j| \\
&< \sum_{j=1}^n \frac{\epsilon}{mn} |v_j| \\
&\le \sum_{j=1}^n \frac{\epsilon}{mn} \|\mathbf{x} - \mathbf{x}_0\| = \frac{\epsilon}{m} \|\mathbf{x} - \mathbf{x}_0\|
\end{align*}

\vspace{0.8em}
\noindent \textbf{Step 4: Extend this to the vector-valued function $f$}\\
\begin{align*}
\|f(\mathbf{x}) - f(\mathbf{x}_0) - L(\mathbf{x} - \mathbf{x}_0)\| &\le \sum_{i=1}^m \left| f_i(\mathbf{x}) - f_i(\mathbf{x}_0) - \sum_{j=1}^n \frac{\partial f_i}{\partial x_j}(\mathbf{x}_0) v_j \right| \\
&< \sum_{i=1}^m \frac{\epsilon}{m} \|\mathbf{x} - \mathbf{x}_0\| = \epsilon \|\mathbf{x} - \mathbf{x}_0\|
\end{align*}

Dividing both sides by $\|\mathbf{x} - \mathbf{x}_0\|$, we get:
$$\frac{\|f(\mathbf{x}) - f(\mathbf{x}_0) - L(\mathbf{x} - \mathbf{x}_0)\|}{\|\mathbf{x} - \mathbf{x}_0\|} < \epsilon$$
Since this holds for any $\epsilon > 0$ when $\|\mathbf{x} - \mathbf{x}_0\|$ is sufficiently small, we conclude:
$$\lim_{\mathbf{x} \to \mathbf{x}_0} \frac{\|f(\mathbf{x}) - f(\mathbf{x}_0) - L(\mathbf{x} - \mathbf{x}_0)\|}{\|\mathbf{x} - \mathbf{x}_0\|} = 0$$
This proves that $f$ is differentiable at $\mathbf{x}_0$, and its derivative is given by $L$.
\end{pf}
\end{dbox}

\begin{lemma}[Chain rule along a coordinate segment]
Let $E \subset \mathbb{R}^n$ and $f: E \to \mathbb{R}^m$. Fix $1 \le i \le m$ and $1 \le j \le n$. Define
$$ \varphi_i^j(t) := f_i(\mathbf{z} + t v \mathbf{e}_j), \quad 0 \le t \le 1, $$
where $\mathbf{z} \in E$ and $v \in \mathbb{R}$. Assume that for some $t_0 \in (0, 1)$ there exists $\rho > 0$ such that
$$ \mathbf{z} + (t_0 + h)v \mathbf{e}_j \in E \quad \text{for all } |h| < \rho, $$
and that the partial derivative $\frac{\partial f_i}{\partial x_j}$ exists at the point $\mathbf{p} := \mathbf{z} + t_0 v \mathbf{e}_j$. Then $\varphi_i^j$ is differentiable at $t_0$ and
$$ (\varphi_i^j)'(t_0) = \frac{\partial f_i}{\partial x_j}(\mathbf{p}) v. $$
\label{coordinate chain rule}
\end{lemma}


\begin{dbox}
\begin{pf}
Let 
\begin{align*}
(\varphi_i^j)'(t_0) = \lim_{h\to 0}\frac{\varphi_i^j(t_0+h)-\varphi_i^j(t_0)}{h}
&=\lim_{h\to 0}\frac{f_i(\mathbf{z} + (t_0+h) v \mathbf{e}_j)-f_i(\mathbf{z} + t_0 v \mathbf{e}_j)}{t}\\
&=\lim_{h\to 0}\frac{f_i(\mathbf{p}+hv\mathbf{e}_j)-f_i(\mathbf{p})}{h}.
\end{align*}
If $v=0$, $f_i(\mathbf{p}+hv\mathbf{e}_j)-\mathbf{p}=0.$
Now consider $v\ne 0$. 

For arbitrarily small $h$, $\mathbf{z} + (t_0 + h)v \mathbf{e}_j=\mathbf{p}+hv\mathbf{e}_j\in B(\mathbf{p},\rho)\subseteq E.$ 
Set $hv=s$, the above quotient becomes:
\[\lim_{s\to 0}\frac{f_i(\mathbf{p}+s\mathbf{e}_j)-f_i(\mathbf{p})}{s}\cdot v = \frac{\partial f_i}{\partial x_j}(\mathbf{p})v,\]
which shows that $\varphi_i^j$ is differentiable at $t_0$ and
$(\varphi_i^j)'(t_0) = \frac{\partial f_i}{\partial x_j}(\mathbf{p}) v.$
\end{pf}
\end{dbox}



\begin{remark}
Let $E \subseteq \mathbb{R}^n$, $f: E \to \mathbb{R}^m$ be a function, and $\mathbf{x}_0$ be an interior point of $E$.
If $f$ is differentiable at $\mathbf{x}_0$, by defintion, there exists a linear map such that the relative error converge to zero near $\mathbf{x}_0.$
Then we call this linear map (total) derivative of $f$ at $\mathbf{x}_0,$ denoted by

\begin{align*}
Df(\mathbf{x}_0) = f^\prime(\mathbf{x}_0) &=
\begin{pmatrix}
- & \nabla f_1(\mathbf{x}_0) & - \\
- & \nabla f_2(\mathbf{x}_0) & - \\
  & \vdots & \\
- & \nabla f_m(\mathbf{x}_0) & - 
\end{pmatrix}_{m \times n}\\ \\ &= 
\begin{pmatrix}
| & | & & | \\
\frac{\partial f}{\partial x_1}(\mathbf{x}_0) & \frac{\partial f}{\partial x_2}(\mathbf{x}_0) & \cdots & \frac{\partial f}{\partial x_n}(\mathbf{x}_0) \\
| & | & & |
\end{pmatrix}_{m \times n}.    
\end{align*}
\end{remark}

Hence, we have:
\begin{enumerate}
    \item The partial derivative is simply the directional derivative along the standard basis vector $\mathbf{e}_j:$
    $$ \frac{\partial f}{\partial x_j}(\mathbf{x}_0) = D_{\mathbf{e}_j} f(\mathbf{x}_0) = f'(\mathbf{x}_0) \mathbf{e}_j $$
    
    \item If $f$ is differentiable, then for any direction vector $\mathbf{v} = \sum_{j=1}^n v_j \mathbf{e}_j$, the directional derivative can be expanded as:
    $$ D_{\mathbf{v}} f(\mathbf{x}_0) = f'(\mathbf{x}_0) \mathbf{v} = f'(\mathbf{x}_0) \left( \sum_{j=1}^n v_j \mathbf{e}_j \right) = \sum_{j=1}^n v_j f'(\mathbf{x}_0) \mathbf{e}_j = \sum_{j=1}^n v_j \frac{\partial f}{\partial x_j}(\mathbf{x}_0) $$
    This can be naturally represented in matrix multiplication form (where the derivative is a $m \times n$ matrix):
    $$ \begin{pmatrix}
    | & | & & | \\
    \frac{\partial f}{\partial x_1}(\mathbf{x}_0) & \frac{\partial f}{\partial x_2}(\mathbf{x}_0) & \cdots & \frac{\partial f}{\partial x_n}(\mathbf{x}_0) \\
    | & | & & |
    \end{pmatrix}_{m \times n}
    \begin{pmatrix}
    v_1 \\
    v_2 \\
    \vdots \\
    v_n
    \end{pmatrix}_{n \times 1}. $$
\end{enumerate}



\begin{definition}[Frobenius Norm]
Let $A=(a_{ij})_{i,j}\in M_{m\times n}(\mathbb{R})$. We can regard $A$ as a vector in $\mathbb{R}^{mn}$, and define its
Frobenius norm by
$$\|A\|_F = \sqrt{\sum_{i=1}^m \sum_{j=1}^n a_{ij}^2}.$$
\end{definition}

\begin{remark}
The Frobenius norm can also be elegantly expressed using the trace of a matrix. 
For $A \in M_{m \times n}(\mathbb{R})$, consider the matrix product $A A^\top$. By the definition of matrix multiplication, the $i$-th diagonal entry of $A A^\top$ is given by:
$$(A A^\top)_{ii} = \sum_{j=1}^n a_{ij} (A^\top)_{ji} = \sum_{j=1}^n a_{ij}^2$$
Since the trace of a matrix is the sum of its diagonal entries, taking the trace of $A A^\top$ yields:
$$\text{tr}(A A^\top) = \sum_{i=1}^m (A A^\top)_{ii} = \sum_{i=1}^m \sum_{j=1}^n a_{ij}^2$$
Therefore, the Frobenius norm can be equivalently written as:
$$\|A\|_F = \sqrt{\text{tr}(A A^\top)}.$$
\end{remark}


\begin{theorem}[Sub-multiplicativity of Frobenius Norm]
Let $A \in M_{m \times n}(\mathbb{R})$ and $B \in M_{n \times p}(\mathbb{R})$. Then the Frobenius norm satisfies the sub-multiplicative property:
$$ \|AB\|_F \le \|A\|_F \|B\|_F $$
In particular, for every column vector $\mathbf{x} \in \mathbb{R}^n$, we have $\|A\mathbf{x}\| \le \|A\|_F \|\mathbf{x}\|$. This implies that the linear map $\mathbf{x} \mapsto A\mathbf{x}$ is continuous on $\mathbb{R}^n$.
\end{theorem}

\begin{dbox}
\begin{pf}
By the definition of matrix multiplication, the $(i, j)$-th entry of the product matrix $AB$ is given by:
$$(AB)_{ij} = \sum_{k=1}^n a_{ik} b_{kj}$$

Applying the Cauchy-Schwarz inequality to the vectors $(a_{i1}, \dots, a_{in})$ and $(b_{1j}, \dots, b_{nj})$, we obtain an upper bound for the square of this entry:
$$|(AB)_{ij}|^2 = \left| \sum_{k=1}^n a_{ik} b_{kj} \right|^2 \le \left( \sum_{k=1}^n a_{ik}^2 \right) \left( \sum_{k=1}^n b_{kj}^2 \right)$$

Now, we compute the squared Frobenius norm of the matrix $AB$ by summing over all entries $(1 \le i \le m$ and $1 \le j \le p)$:
\begin{align*}
\|AB\|_F^2 &= \sum_{i=1}^m \sum_{j=1}^p (AB)_{ij}^2 \\
&\le \sum_{i=1}^m \sum_{j=1}^p \left( \sum_{k=1}^n a_{ik}^2 \right) \left( \sum_{k=1}^n b_{kj}^2 \right)
\end{align*}

Notice that the sum over $i$ only depends on the terms $a_{ik}$, and the sum over $j$ only depends on the terms $b_{kj}$. Therefore, we can separate and factor the summations:
\begin{align*}
\|AB\|_F^2 &\le \left( \sum_{i=1}^m \sum_{k=1}^n a_{ik}^2 \right) \left( \sum_{j=1}^p \sum_{k=1}^n b_{kj}^2 \right) \\
&= \|A\|_F^2 \cdot \|B\|_F^2
\end{align*}

Finally, taking the square root of both sides, we get:
$$ \|AB\|_F \le \|A\|_F \|B\|_F $$
This completes the proof.
\end{pf}
\end{dbox}

\clearpage

\begin{theorem}[Chain Rule]
Let $E \subseteq \mathbb{R}^n,\, F\subseteq\mathbb{R}^m,\, f: E \to F,\, g: F\to\mathbb{R}^p$ and $\mathbb{x}_0\in int(E).$ 
Suppose $f$ is differentiable at $\mathbf{x}_0$, $f(\mathbf{x}_0)$ lies in the interior of $E$, and $g$ is differentiable at $f(\mathbf{x}_0)$.
Then $g \circ f$ is differentiable at $\mathbf{x}_0$, and
$$(g \circ f)'(\mathbf{x}_0) = g'(f(\mathbf{x}_0)) f'(\mathbf{x}_0).$$
\end{theorem}

\begin{dbox}
\begin{pf}
Let $\mathbf{y}_0 = f(\mathbf{x}_0)$ and denote the derivatives by $A = f'(\mathbf{x}_0): \mathbb{R}^n \to \mathbb{R}^m$ and $B = g'(f(\mathbf{x}_0)) = g'(\mathbf{y}_0): \mathbb{R}^m \to \mathbb{R}^p$.

\vspace{0.8em}
\noindent \textbf{Step 1: Define the error function $\eta$} \\
Define $\eta: \mathbb{R}^m \to \mathbb{R}^p$ by:
$$
\eta(\mathbf{k}) = 
\begin{cases}
\frac{g(\mathbf{y}_0+\mathbf{k}) - g(\mathbf{y}_0) - B\mathbf{k}}{\|\mathbf{k}\|} & \text{if } \mathbf{k} \ne \mathbf{0} \\
\mathbf{0} & \text{if } \mathbf{k} = \mathbf{0}
\end{cases}
$$ 

Because $g$ is differentiable at $\mathbf{y}_0$, $\lim_{\mathbf{k} \to \mathbf{0}} \eta(\mathbf{k}) = \mathbf{0}$, which implies $\eta$ is continuous at $\mathbf{0}$.
Moreover, the relation
$$g(\mathbf{y}_0+\mathbf{k}) - g(\mathbf{y}_0) - B\mathbf{k} = \eta(\mathbf{k}) \|\mathbf{k}\|$$
is true for all $\mathbf{k} \in \mathbb{R}^m$.

\vspace{0.8em}
\noindent \textbf{Step 2: Expand the composition difference} \\
Take $\mathbf{h}$ such that $\mathbf{x}_0+\mathbf{h} \in E$. Let $\mathbf{k} = f(\mathbf{x}_0+\mathbf{h}) - f(\mathbf{x}_0)$.
Consider the difference:
\begin{align*}
(g \circ f)(\mathbf{x}_0+\mathbf{h}) - (g \circ f)(\mathbf{x}_0) 
&= g(f(\mathbf{x}_0+\mathbf{h})) - g(f(\mathbf{x}_0)) \\
&= g(\mathbf{y}_0+\mathbf{k}) - g(\mathbf{y}_0) \\
&= B(f(\mathbf{x}_0+\mathbf{h}) - f(\mathbf{x}_0)) \\
&\quad + \eta(f(\mathbf{x}_0+\mathbf{h}) - f(\mathbf{x}_0)) \|f(\mathbf{x}_0+\mathbf{h}) - f(\mathbf{x}_0)\|
\end{align*}

Subtracting $BA\mathbf{h}$ from both sides, we get:
\begin{align*}
&(g \circ f)(\mathbf{x}_0+\mathbf{h}) - (g \circ f)(\mathbf{x}_0) - BA\mathbf{h} \\
&= B(f(\mathbf{x}_0+\mathbf{h}) - f(\mathbf{x}_0) - A\mathbf{h}) + \eta(f(\mathbf{x}_0+\mathbf{h}) - f(\mathbf{x}_0)) \|f(\mathbf{x}_0+\mathbf{h}) - f(\mathbf{x}_0)\|
\end{align*} 

\vspace{0.8em}
\noindent \textbf{Step 3: Estimate the first term} \\
Note that:
$$\frac{\|B(f(\mathbf{x}_0+\mathbf{h}) - f(\mathbf{x}_0) - A\mathbf{h})\|}{\|\mathbf{h}\|} \le \|B\| \frac{\|f(\mathbf{x}_0+\mathbf{h}) - f(\mathbf{x}_0) - A\mathbf{h}\|}{\|\mathbf{h}\|}$$ 

By the assumption that $f$ is differentiable at $\mathbf{x}_0$, $\lim_{\mathbf{h} \to \mathbf{0}} \frac{\|f(\mathbf{x}_0+\mathbf{h}) - f(\mathbf{x}_0) - A\mathbf{h}\|}{\|\mathbf{h}\|} = 0$. By the Squeeze Theorem, the first term converges to zero.

\vspace{0.8em}
\noindent \textbf{Step 4: Estimate the second term} \\
Since $f$ is differentiable at $\mathbf{x}_0$, $f$ is continuous at $\mathbf{x}_0$, meaning $\lim_{\mathbf{h} \to \mathbf{0}} (f(\mathbf{x}_0+\mathbf{h}) - f(\mathbf{x}_0)) = \mathbf{0}$.
On the other hand, $\eta$ is continuous at $\mathbf{0}$ (and $\lim_{\mathbf{k} \to \mathbf{0}} \eta(\mathbf{k}) = \mathbf{0}$), it follows that $\lim_{\mathbf{h} \to \mathbf{0}} \eta(f(\mathbf{x}_0+\mathbf{h}) - f(\mathbf{x}_0)) = \mathbf{0}$

Next, consider:
\begin{align*}
\frac{\|f(\mathbf{x}_0+\mathbf{h}) - f(\mathbf{x}_0)\|}{\|\mathbf{h}\|} &= \frac{\|f(\mathbf{x}_0+\mathbf{h}) - f(\mathbf{x}_0) - A\mathbf{h} + A\mathbf{h}\|}{\|\mathbf{h}\|} \\
&\le \frac{\|f(\mathbf{x}_0+\mathbf{h}) - f(\mathbf{x}_0) - A\mathbf{h}\|}{\|\mathbf{h}\|} + \frac{\|A\mathbf{h}\|}{\|\mathbf{h}\|} \\
&\le \underbrace{\frac{\|f(\mathbf{x}_0+\mathbf{h}) - f(\mathbf{x}_0) - A\mathbf{h}\|}{\|\mathbf{h}\|}}_{\to 0\text{ by differentiability}} + \underbrace{\|A\|}_{fixed}
\end{align*} 

By the Squeeze Theorem:
$$\lim_{\mathbf{h} \to \mathbf{0}} \frac{\|\eta(f(\mathbf{x}_0+\mathbf{h}) - f(\mathbf{x}_0))\| \|f(\mathbf{x}_0+\mathbf{h}) - f(\mathbf{x}_0)\|}{\|\mathbf{h}\|} = 0$$

Combining Step 3 and Step 4, we conclude:
$$\lim_{\mathbf{h} \to \mathbf{0}} \frac{\|(g \circ f)(\mathbf{x}_0+\mathbf{h}) - (g \circ f)(\mathbf{x}_0) - BA\mathbf{h}\|}{\|\mathbf{h}\|} = 0$$ 
This proves that $g \circ f$ is differentiable at $\mathbf{x}_0$ with derivative $BA = g'(f(\mathbf{x}_0)) f'(\mathbf{x}_0)$.
\end{pf}
\end{dbox}

\begin{definition}
    Let $E\subseteq\mathbf{R}^n$ be open and $f:E\to\mathbf{R}^m.$
    We say $f$ is twice contiuously differentiable on $E$, denoted by $f\in C^2(E)$.
    \begin{enumerate}
    \item $f$ is continuously differentiable (its first order partial derivatives exist and are continuous).
    \item For each $i=1,2,\cdots n,$ the partial derivative $\frac{\partial f}{\partial x_i}: E\to\mathbf{R}^m$ is continously differentiable on $E$ (its second order partial derivatives exist and are continous on $E$).
\end{enumerate}
\end{definition}


\begin{theorem}[Clairaut]
Let $E$ be an open set and $f \in C^2(E)$. Then for all $1 \le i, j \le n$ and all $\mathbf{x}_0 \in E$:
$$\frac{\partial^2 f}{\partial x_i \partial x_j}(\mathbf{x}_0) = \frac{\partial^2 f}{\partial x_j \partial x_i}(\mathbf{x}_0).$$
\end{theorem}

\begin{dbox}
\begin{pf}
Let $f = (f_1, f_2, \dots, f_m)$ where $f_k: E \to \mathbb{R}$. Without loss of generality, we can assume $f: E \to \mathbb{R}$ and the point of interest is $\mathbf{x}_0 = \mathbf{0}$. (By some shifting, $F(\mathbf{x}) = f(\mathbf{x} + \mathbf{x}_0) \Rightarrow F(\mathbf{0}) = f(\mathbf{x}_0)$). We also assume $i \neq j$.

Since $f \in C^2(E)$ and $\mathbf{0} \in E$, given any $\epsilon > 0$, there exists $\delta_0 > 0$ such that:
$$ \left| \frac{\partial^2 f}{\partial x_j \partial x_i}(\mathbf{x}) - \frac{\partial^2 f}{\partial x_j \partial x_i}(\mathbf{0}) \right| \le \frac{\epsilon}{2} \quad \text{and} \quad \left| \frac{\partial^2 f}{\partial x_i \partial x_j}(\mathbf{x}) - \frac{\partial^2 f}{\partial x_i \partial x_j}(\mathbf{0}) \right| \le \frac{\epsilon}{2} $$
if $\|\mathbf{x}\| < 2\delta_0$.

\vspace{0.8em}
\noindent \textbf{Step 1: The Second-Order Difference (DiD Approximation)}\\
Fix $\delta < \delta_0$, and consider the rectangle. Define the second-order difference $X$:
$$ X = f(\delta \mathbf{e}_i + \delta \mathbf{e}_j) - f(\delta \mathbf{e}_i) - f(\delta \mathbf{e}_j) + f(\mathbf{0}) $$
We can group the terms in two different ways:
\begin{align}
X &= \left( f(\delta \mathbf{e}_i + \delta \mathbf{e}_j) - f(\delta \mathbf{e}_i) \right) - \left( f(\delta \mathbf{e}_j) - f(\mathbf{0}) \right) \tag{1} \\
X &= \left( f(\delta \mathbf{e}_i + \delta \mathbf{e}_j) - f(\delta \mathbf{e}_j) \right) - \left( f(\delta \mathbf{e}_i) - f(\mathbf{0}) \right) \tag{2}
\end{align}

\vspace{0.8em}
\noindent \textbf{Step 2: Apply Mean Value Theorem for Grouping (1)}\\
Let $\phi(t) = f(t \mathbf{e}_i + \delta \mathbf{e}_j) - f(t \mathbf{e}_i)$. Then from (1), we have:
$$ X = \phi(\delta) - \phi(0) $$
By the MVT, $X = \delta \phi'(c_i)$ for some $c_i \in (0, \delta)$ \quad (*).

Note that $\phi'(t) = \frac{\partial f}{\partial x_i}(t \mathbf{e}_i + \delta \mathbf{e}_j) - \frac{\partial f}{\partial x_i}(t \mathbf{e}_i)$.
So, $\phi'(c_i) = \frac{\partial f}{\partial x_i}(c_i \mathbf{e}_i + \delta \mathbf{e}_j) - \frac{\partial f}{\partial x_i}(c_i \mathbf{e}_i)$.

Now, let $g(s) = \frac{\partial f}{\partial x_i}(c_i \mathbf{e}_i + s \mathbf{e}_j)$. Then $g'(s) = \frac{\partial^2 f}{\partial x_j \partial x_i}(c_i \mathbf{e}_i + s \mathbf{e}_j)$.
We can rewrite (*) as:
$$ \phi'(c_i) = g(\delta) - g(0) $$
By applying MVT again to $g(s)$, there exists $c_j \in (0, \delta)$ such that $g(\delta) - g(0) = \delta g'(c_j)$.
Thus, $\phi'(c_i) = \delta \frac{\partial^2 f}{\partial x_j \partial x_i}(c_i \mathbf{e}_i + c_j \mathbf{e}_j)$ \quad (**).

From (*) and (**), we obtain:
$$ X = \delta^2 \frac{\partial^2 f}{\partial x_j \partial x_i}(c_i \mathbf{e}_i + c_j \mathbf{e}_j) \quad \text{where } 0 < c_i < \delta, \ 0 < c_j < \delta. $$

\vspace{0.8em}
\noindent \textbf{Step 3: Symmetric Argument for Grouping (2)}\\
Similarly, starting from (2) and taking derivatives first with respect to $x_j$ and then $x_i$, MVT gives us:
$$ X = \delta^2 \frac{\partial^2 f}{\partial x_i \partial x_j}(d_i \mathbf{e}_i + d_j \mathbf{e}_j) \quad \text{where } 0 < d_i < \delta, \ 0 < d_j < \delta. $$

\vspace{0.8em}
\noindent \textbf{Step 4: Bounding the Error using Continuity}\\
Notice that the norm of the points we found is strictly bounded:
$$ \|c_i \mathbf{e}_i + c_j \mathbf{e}_j\| = \sqrt{c_i^2 + c_j^2} < \sqrt{\delta^2 + \delta^2} = \sqrt{2}\delta < \sqrt{2}\delta_0 < 2\delta_0 $$
Because these points fall within the $2\delta_0$ radius, our initial continuity conditions apply:
$$ \left| \frac{\partial^2 f}{\partial x_j \partial x_i}(c_i \mathbf{e}_i + c_j \mathbf{e}_j) - \frac{\partial^2 f}{\partial x_j \partial x_i}(\mathbf{0}) \right| \le \epsilon \implies \left| \frac{X}{\delta^2} - \frac{\partial^2 f}{\partial x_j \partial x_i}(\mathbf{0}) \right| \le  \frac{\epsilon}{2} $$
Similarly,
$$ \left| \frac{X}{\delta^2} - \frac{\partial^2 f}{\partial x_i \partial x_j}(\mathbf{0}) \right| \le  \frac{\epsilon}{2} $$

By the triangle inequality:
$$ \left| \frac{\partial^2 f}{\partial x_i \partial x_j}(\mathbf{0}) - \frac{\partial^2 f}{\partial x_j \partial x_i}(\mathbf{0}) \right| \le \epsilon \quad \forall \epsilon > 0 $$
Hence, we conclude that:
$$ \frac{\partial^2 f}{\partial x_i \partial x_j}(\mathbf{0}) = \frac{\partial^2 f}{\partial x_j \partial x_i}(\mathbf{0}). $$
\end{pf}
\end{dbox}

\begin{rmk}[Numerical Approximation via Second-Order Difference]
From the proof above, we know that by the continuity of $\frac{\partial^2 f}{\partial x_j \partial x_i}$ and $\frac{\partial^2 f}{\partial x_i \partial x_j}$,
$$ \lim_{\delta \to 0} \frac{X}{\delta^2} = \frac{\partial^2 f}{\partial x_j \partial x_i}(\mathbf{0}) \quad \text{and} \quad \lim_{\delta \to 0} \frac{X}{\delta^2} = \frac{\partial^2 f}{\partial x_i \partial x_j}(\mathbf{0}) $$
Hence, if $f \in C^2(E)$, for any point $\mathbf{x}_0$, we can use numerical approximation by a second-order finite difference without computing the derivative directly:
$$ \lim_{\delta \to 0} \frac{f(\mathbf{x}_0 + \delta \mathbf{e}_i + \delta \mathbf{e}_j) - f(\mathbf{x}_0 + \delta \mathbf{e}_i) - f(\mathbf{x}_0 + \delta \mathbf{e}_j) + f(\mathbf{x}_0)}{\delta^2} = \frac{\partial^2 f}{\partial x_j \partial x_i}(\mathbf{x}_0) = \frac{\partial^2 f}{\partial x_i \partial x_j}(\mathbf{x}_0) $$
(In economics, this logic is highly analogous to the DiD).
\end{rmk}




















\end{document}